%% =============================================================
%%  Capitolo 2 — Fondamenti
%%  Sistemi a regole, retrieval-augmented generation
%%  e approcci neuro-simbolici
%%
%%  Tesi di Laurea Triennale — Francesco Zamar
%%  Università degli Studi di Trieste, DIA
%%  Relatore: Prof. Andrea De Lorenzo
%%  A.A. 2025/2026
%% =============================================================

\chapter{Fondamenti}
\label{ch:background}

Il \Cref{ch:contesto} ha mostrato perché la verifica di una Nota AIFA (Agenzia Italiana del Farmaco) è un
problema clinico e regolatorio insieme, e perché un CDSS (\emph{Clinical Decision Support System}) rule-based puro,
pur essendo lo stato dell'arte nei gestionali italiani, non basta a produrre
quella spiegazione tracciabile che il medico richiede di fronte a una decisione
che giudica controintuitiva. Per costruire un sistema che affronti entrambi i
limiti servono tre strumenti tecnici, di cui questo capitolo offre un'esposizione
mirata al solo livello di dettaglio che i capitoli successivi assumeranno come
acquisito. Il primo è un linguaggio per rappresentare l'incertezza clinica
quando un parametro di laboratorio è assente: la logica trivalente di Kleene,
discussa nella \Cref{sec:rule-based} insieme al paradigma dei sistemi a regole
che la ospita. Il secondo è il meccanismo con cui ancorare le spiegazioni al
testo normativo ufficiale invece che alla memoria parametrica di un modello
generativo: la \Cref{sec:rag} descrive l'architettura RAG e le sue varianti di
retrieval. Il terzo è il framework concettuale che permette di combinare i due:
la \Cref{sec:neurosymbolic} colloca il sistema nel paesaggio degli approcci
neuro-simbolici e motiva la scelta dell'accoppiamento debole. Il capitolo si
chiude con il vocabolario delle metriche che ricompariranno nel
\Cref{ch:valutazione}.


% ------------------------------------------------------------------
\section{Sistemi a regole e reasoning simbolico}
\label{sec:rule-based}
% ------------------------------------------------------------------

La nozione di ragionamento automatico su base simbolica ha radici profonde
nell'informatica e nella logica matematica. Un sistema a regole, detto
anche \emph{sistema di produzione} o \emph{rule engine}, rappresenta la
conoscenza come un insieme di proposizioni condizionali della forma
\textit{se condizione allora azione}, che un motore di inferenza valuta
sistematicamente su un insieme di fatti~\cite{hripcsak1994arden}. La
semplicità della struttura nasconde una potenza espressiva considerevole: ogni
regola è leggibile da un esperto di dominio, ogni passo inferenziale è
tracciabile e ogni decisione può essere motivata elencando le regole che
hanno contribuito ad attivarla. Queste proprietà rendono i sistemi a regole
particolarmente adatti a contesti safety-critical, dove la trasparenza del
processo decisionale non è un lusso ma un requisito
normativo~\cite{greenes2007cdss}.

Il meccanismo di inferenza più comune nei sistemi a regole orientati ai fatti
è il \emph{forward chaining}: partendo dai fatti noti, il motore cerca le
regole le cui condizioni siano soddisfatte, ne applica le azioni, aggiunge i
nuovi fatti all'insieme di lavoro e itera fino a quando non è più possibile
attivare alcuna regola. L'alternativa, il \emph{backward chaining}, parte
dall'obiettivo da dimostrare e cerca a ritroso le condizioni che lo
implicano; è la strategia tipica dei sistemi Prolog-like e degli oracoli
di interrogazione. Il forward chaining è più naturale per la verifica di
ammissibilità normativa, dove si parte da un insieme di parametri clinici
noti e si vuole determinare se soddisfano cumulativamente i criteri di una
Nota AIFA: il motore applica le regole nell'ordine in cui i fatti
si rendono disponibili, senza bisogno di conoscere a priori l'obiettivo.

Una proprietà del forward chaining rilevante per l'applicazione clinica è
la \emph{monotonicità} dell'inferenza rispetto all'acquisizione di nuovi
dati: aggiungere un fatto all'insieme di lavoro non può rendere false le
conclusioni già raggiunte a meno che non si disponga di regole di
retrazione esplicita. Questa caratteristica, nota anche come \emph{closed-world
monotonicity} nel senso in cui l'insieme dei fatti cresce ma non contrae
le conclusioni già validate, è importante perché garantisce che la
decisione prodotta dal rule engine rimanga stabile durante la sessione
clinica, finché non arrivano nuovi dati contraddittori.

% ──────────────────────────────────────────────────────
\subsection{Logica di Kleene a tre valori}
\label{subsec:kleene}
% ──────────────────────────────────────────────────────

La logica classica binaria postula che ogni proposizione sia vera o falsa.
In un contesto clinico reale, questo assunto è quasi sempre violato: un
parametro di laboratorio può essere non ancora disponibile, un referto può
essere atteso, un campo della cartella clinica può non essere stato compilato.
Un sistema binario che gestisce queste situazioni come \emph{false} è
strutturalmente inadeguato, perché equipara l'assenza di informazione alla
confutazione. La logica trivalente di Kleene (K3)~\cite{kleene1952introduction}
introduce un terzo valore, convenzionalmente indicato come $\mathbf{U}$
(\emph{Undefined} o \emph{Unknown}), che denota una proposizione il cui valore
di verità non è ancora determinabile. Bochvar aveva formulato
indipendentemente un calcolo analogo~\cite{bochvar1938}, ma è la
caratterizzazione algebrica di Kleene a essere adottata nei sistemi a regole
moderni per la sua coerenza con le operazioni proposizionali standard.

Formalmente, sia $\mathcal{V} = \{\mathbf{T}, \mathbf{F}, \mathbf{U}\}$
l'insieme dei valori di verità. Le operazioni fondamentali si definiscono
come estensione delle tavole di verità classiche con il principio di
minima informazione: se almeno un operando è $\mathbf{U}$ e il risultato
non è determinabile dagli altri operandi, il risultato è $\mathbf{U}$.
Le tavole di verità per congiunzione e disgiunzione rispettano questo principio:
\[
  A \land B = \min(A,\, B) \quad\text{e}\quad A \lor B = \max(A,\, B)
\]
rispetto all'ordine $\mathbf{F} < \mathbf{U} < \mathbf{T}$. La negazione
si comporta classicamente su $\mathbf{T}$ e $\mathbf{F}$, mentre
$\lnot\mathbf{U} = \mathbf{U}$. Alcune proprietà notevoli emergono
immediatamente: $A \land \mathbf{F} = \mathbf{F}$ indipendentemente da $A$,
il che significa che una condizione falsificante nota è sufficiente a
determinare un risultato negativo anche in presenza di dati mancanti; dualmente,
$A \lor \mathbf{T} = \mathbf{T}$ indipendentemente da $A$.

Oltre all'ordine di verità, K3 ammette un distinto \emph{ordine informativo}
$\sqsubseteq$ definito da $\mathbf{U} \sqsubseteq \mathbf{T}$ e
$\mathbf{U} \sqsubseteq \mathbf{F}$, con $\mathbf{T}$ e $\mathbf{F}$
incomparabili tra loro. Rispetto a $\sqsubseteq$, le operazioni $\land$,
$\lor$ e $\lnot$ sono monotone: se $A \sqsubseteq A'$ e $B \sqsubseteq B'$,
allora $A \land B \sqsubseteq A' \land B'$ e $A \lor B \sqsubseteq A' \lor B'$.
Operativamente questo significa che acquisire nuovi dati clinici, cioè
sostituire $\mathbf{U}$ con un valore determinato, raffina monotonicamente
l'esito della formula complessiva, senza mai ribaltarne arbitrariamente il
segno: il sistema può passare da $\mathbf{U}$ a $\mathbf{T}$ o a $\mathbf{F}$,
ma sempre in modo coerente con la struttura logica delle regole.

Nel sistema di questa tesi, il valore $\mathbf{U}$ corrisponde al campo clinico
non compilato o al parametro non ancora misurato. Quando tutte le condizioni
di una Nota AIFA producono $\mathbf{T}$, il rule engine emette la decisione
\decision{RIMBORSABILE}; quando almeno una condizione produce $\mathbf{F}$
e nessun cammino alternativo è percorribile, emette \decision{NON\_RIMBORSABILE};
quando la presenza di $\mathbf{U}$ in posizione critica rende il risultato
non determinabile con i dati disponibili, emette \decision{NON\_DETERMINABILE}.
Questa tricotomia è un'istanza diretta della semantica K3 e la sua
implementazione nel motore è verificata analiticamente da test di unità
sulle tavole di verità.

Il quadro storico dei sistemi esperti a regole in medicina (da MYCIN
\cite{shortliffe1974mycin} ad Arden Syntax \cite{hripcsak1994arden} ai CDSS
moderni \cite{greenes2007cdss,berner2007cdss,sutton2020overview}) è già
stato richiamato nella \Cref{sec:cdss-overview}. Qui basta osservare che i
sistemi rule-based puri offrono tracciabilità completa e aggiornabilità
normativa ma soffrono di rigidità espressiva, mentre i modelli di
apprendimento automatico offrono espressività ma soffrono di opacità: nessuna
delle due famiglie è da sola sufficiente per un sistema che debba essere allo
stesso tempo corretto per costruzione sul piano normativo e capace di produrre
una spiegazione in linguaggio naturale auditabile. È la motivazione che porta,
nella \Cref{sec:neurosymbolic}, alla scelta neuro-simbolica.


% ------------------------------------------------------------------
\section{Recupero e generazione condizionata}
\label{sec:rag}
% ------------------------------------------------------------------

I grandi modelli linguistici (\emph{Large Language Model}, LLM) sono addestrati su corpora testuali di
dimensioni enormi e acquisiscono conoscenza del mondo in forma distribuita
nei loro parametri. Questa \emph{conoscenza parametrica} è però statica:
il modello non può aggiornare autonomamente ciò che sa senza un nuovo ciclo
di addestramento, e non ha accesso diretto a documenti che non ha visto
durante il training. In un dominio a conoscenza chiusa come le Note AIFA,
dove i criteri di rimborsabilità cambiano per decreto ministeriale e devono
essere applicati nella loro formulazione letterale, fare affidamento
esclusivo sulla memoria parametrica del modello produce due categorie di
errori: l'allucinazione di criteri inesistenti e l'omissione di criteri
aggiornati dopo la data di cutoff~\cite{ji2023hallucination}. La
\emph{Retrieval-Augmented Generation} (RAG) è stata proposta da Lewis et al.
come soluzione a questo problema: invece di rispondere esclusivamente da
memoria, il modello riceve come contesto i passaggi documentali recuperati
da un corpus esterno che il sistema può aggiornare indipendentemente
dall'addestramento~\cite{lewis2020rag}.

% ──────────────────────────────────────────────────────
\subsection{Architettura di un sistema RAG}
\label{subsec:rag-arch}
% ──────────────────────────────────────────────────────

Una pipeline RAG canonica si articola in due fasi temporalmente separate.
La fase \emph{offline}, eseguita una tantum o ad ogni aggiornamento
del corpus, costruisce un indice dei documenti: il testo viene segmentato
in unità, ciascuna unità viene convertita in un vettore numerico tramite un
modello di embedding, e i vettori vengono memorizzati in un database
specializzato per la ricerca per similarità. La fase \emph{online}, eseguita
ad ogni interrogazione, converte la domanda dello stesso modello di
embedding in un vettore query, recupera i vettori di documento più simili,
e fornisce al modello generativo tanto la domanda originale quanto i passaggi
recuperati come contesto esplicito. Il modello produce così una risposta
che è condizionata sui documenti, non soltanto sulla propria memoria
parametrica~\cite{gao2023ragsurvey}.

Un sistema RAG combina quindi un componente di recupero, che seleziona dai
documenti i passaggi rilevanti per una domanda, con un modello generativo,
che produce una risposta condizionata su quei passaggi. Quando i documenti
hanno valore normativo (come accade per le Nota AIFA, dove ogni frase
è vincolante) questa architettura riduce gli errori di allucinazione e,
soprattutto, fornisce una catena di citazioni che il revisore umano può
controllare in tempo lineare rispetto al numero di fonti citate.

% ──────────────────────────────────────────────────────
\subsection{Chunking e indicizzazione}
\label{subsec:chunking}
% ──────────────────────────────────────────────────────

Il \emph{chunking}, la segmentazione del testo in unità di dimensione
adatta all'embedding, è una delle scelte progettuali più critiche in
una pipeline RAG, perché condiziona sia la qualità del retrieval sia la
lunghezza del contesto consegnato al modello generativo. Unità troppo grandi
contengono più informazione ma diluiscono il segnale semantico e rischiano
di superare la finestra di contesto del modello; unità troppo piccole sono
più precise ma possono separare passaggi che il modello deve leggere insieme
per costruire una risposta coerente.

Le strategie di chunking si suddividono grossolanamente in due famiglie.
Il \emph{fixed-size chunking} divide il testo ogni $w$ token, eventualmente
con una finestra di sovrapposizione di ampiezza $s < w$ per preservare la
continuità semantica ai confini: produce chunk di dimensione prevedibile e
si implementa in modo semplice, ma ignora la struttura del documento.
Il \emph{semantic chunking} invece sfrutta la struttura del documento
(paragrafi, sezioni, elenchi) per produrre unità semanticamente coerenti;
richiede un analizzatore sintattico o una conoscenza della struttura del
documento, ma produce chunk che corrispondono a unità informative discrete,
particolarmente utili con testi normativi dove ogni paragrafo di solito
esprime un criterio autonomo.

L'indicizzazione dei chunk avviene in un \emph{database vettoriale}, una
struttura dati specializzata per la ricerca per similarità in spazi ad
alta dimensione. I database vettoriali moderni implementano algoritmi di
indicizzazione approssimata come HNSW (\emph{Hierarchical Navigable Small
World})~\cite{chromadb}, che costruiscono un grafo navigabile a più livelli
gerarchici e permettono di trovare i vicini approssimati di un vettore query
in tempo logaritmico rispetto alla dimensione del corpus, a fronte di un
costo di costruzione dell'indice proporzionale al numero di chunk.

% ──────────────────────────────────────────────────────
\subsection{Embedding, similarità densa e reranking}
\label{subsec:embedding}
% ──────────────────────────────────────────────────────

Il modello di embedding è il componente che trasforma testo in vettori.
I modelli basati su trasformatori a doppio encoder (\emph{bi-encoder})
producono una rappresentazione densa per ogni sequenza di input calcolando
l'embedding da una singola passata attraverso il modello~\cite{vaswani2017attention}.
L'architettura originale del trasformatore, introdotta da Vaswani et al.
con il meccanismo di \emph{self-attention}~\cite{vaswani2017attention}, ha
reso possibile la costruzione di encoder contestuali di alta qualità: il
vettore di embedding di un token non è fisso ma dipende da tutti gli altri
token della sequenza, permettendo di catturare relazioni semantiche a lunga
distanza. BERT~\cite{devlin2019bert} ha poi mostrato come il pre-addestramento
su grandi corpora tramite predizione mascherata produca rappresentazioni
trasferibili che, con un minimo di fine-tuning, raggiungono prestazioni
all'avanguardia su decine di task di elaborazione del linguaggio naturale (\emph{Natural Language Processing}, NLP).

Sentence-BERT (SBERT)~\cite{reimers2019sbert} ha risolto il problema
specifico della similarità semantica fra frasi: BERT originale produce
embedding per token, non per sequenze, e la media o l'aggregazione dei token
embedding produce rappresentazioni di scarsa qualità per il confronto
semantico inter-frase. SBERT addestra reti siamesi con pooling sul token
\texttt{[CLS]} o con mean-pooling sull'output dell'encoder, ottimizzando
esplicitamente la similarità coseno fra coppie di frasi simili e dissimili.
Il modello \embedmodel{}, adottato in questo lavoro per l'indicizzazione
dei chunk, è la variante multilingue di SBERT addestrata su 50 lingue,
che garantisce la coerenza delle rappresentazioni fra il testo normativo
italiano e le domande cliniche degli utenti.

La similarità fra un vettore query $\mathbf{q}$ e un vettore chunk
$\mathbf{d}$ si misura con la \emph{cosine similarity}:
\[
  \text{cos}(\mathbf{q}, \mathbf{d}) = \frac{\mathbf{q} \cdot \mathbf{d}}{\|\mathbf{q}\|\,\|\mathbf{d}\|}
\]
che varia tra $-1$ e $1$ ed è invariante per riscalamento dei vettori. Quando
i vettori sono già normalizzati a norma unitaria, come accade nell'implementazione
standard di SBERT e dei database vettoriali, la cosine similarity si riduce
al prodotto scalare $\mathbf{q} \cdot \mathbf{d}$, che può essere calcolato
efficientemente con operazioni BLAS e permette la ricerca del nearest neighbor
approssimato mediante strutture HNSW.

Il retrieval denso tramite bi-encoder è rapido ma approssimato: il bi-encoder
codifica query e documento indipendentemente, senza modellare l'interazione
fra i due. Un approccio alternativo, storicamente dominante prima dell'era
degli embedding neurali, è il \emph{retrieval lessicale} basato sulla
frequenza dei termini. Il modello BM25 (Best Match 25, nella tradizione
del probabilistic relevance framework di Robertson e Zaragoza) assegna ad
ogni documento un punteggio che tiene conto della frequenza del termine nella
query all'interno del documento (\emph{term frequency}, TF) e della rarità
del termine nel corpus (\emph{inverse document frequency}, IDF), con una
funzione di saturazione che attenua l'effetto dei documenti molto lunghi.
Il punteggio BM25 per un documento $d$ rispetto a una query con termine $t$
è:
\[
  \text{BM25}(d,t) = \text{IDF}(t) \cdot
    \frac{\text{TF}(t,d) \cdot (k_1+1)}{\text{TF}(t,d) + k_1 \cdot \left(1 - b + b \cdot \frac{|d|}{\overline{|d|}}\right)}
\]
dove $k_1$ e $b$ sono iperparametri che controllano rispettivamente la
saturazione della frequenza e la normalizzazione per lunghezza, e
$\overline{|d|}$ è la lunghezza media dei documenti nel corpus. BM25 eccelle
quando la query contiene termini tecnici rari e specifici, come i codici
ATC (\emph{Anatomical Therapeutic Chemical}) o i nomi di molecole nelle Note AIFA, ma fatica con le parafrasi
semantiche, dove il documento risponde alla domanda ma usa vocabolario
diverso. Il retrieval denso tramite embedding ha il comportamento opposto:
cattura la vicinanza semantica ma può perdere l'ancoraggio lessicale esatto.

Questa complementarietà ha portato allo sviluppo di pipeline di retrieval
ibride che combinano il punteggio BM25 e il punteggio coseno, e di
un secondo stadio di ordinamento a latenza maggiore ma precisione più alta:
il \emph{cross-encoder reranker}~\cite{nogueira2019passage}. Mentre il
bi-encoder encode query e documento separatamente, il cross-encoder li
concatena come singolo input al modello, permettendo al meccanismo di
self-attention di modellare direttamente le interazioni fra i token della
query e quelli del chunk. Il risultato è uno score di rilevanza più
preciso, a costo di $O(k)$ passate forward aggiuntive per i $k$ candidati
da riordinare. La pipeline \emph{two-stage} (bi-encoder a bassa latenza
per il recupero di un sovrainsieme di $K$ candidati, seguito da cross-encoder
per la selezione finale dei top-$k$) è diventata l'architettura standard
per i sistemi RAG ad alta precisione~\cite{gao2023ragsurvey}. In questo
lavoro, il cross-encoder \reranker{} è addestrato in italiano su coppie di
similarità semantica e opera come secondo stadio di selezione sui chunk
preliminarmente recuperati dal database vettoriale.


% ------------------------------------------------------------------
\section{Approcci neuro-simbolici}
\label{sec:neurosymbolic}
% ------------------------------------------------------------------

Un sistema \emph{neuro-simbolico} integra componenti di apprendimento
automatico (tipicamente reti neurali profonde) con componenti di
ragionamento simbolico (tipicamente regole, ontologie o logiche formali).
La motivazione è duplice: i sistemi puramente simbolici sono precisi e
auditabili ma fragili di fronte a dati rumorosi o incompleti e incapaci di
gestire input non strutturati come il testo libero; i sistemi puramente
neurali generalizzano bene e gestiscono l'incertezza, ma producono ragionamenti
opachi e sono inclini ad allucinazioni su fatti specifici di dominio. Un
sistema neuro-simbolico mira a ereditare i pregi di entrambi, e la sua
realizzazione concreta dipende dal tipo e dal grado di accoppiamento fra le
due componenti~\cite{yu2023neuro}.

La letteratura distingue tradizionalmente fra accoppiamento debole
(\emph{loose coupling}) e accoppiamento stretto (\emph{tight coupling}).
Nell'accoppiamento debole, i due sistemi interagiscono tramite interfacce
ben definite: il modulo simbolico riceve l'output del modulo neurale (ad
esempio un'entità estratta da testo) e vi applica regole; oppure il modulo
neurale riceve l'output del modulo simbolico (ad esempio un fact verificato)
e lo usa come contesto aggiuntivo. I due moduli sono addestrabili e
verificabili in isolamento. Nell'accoppiamento stretto, la componente neurale
e quella simbolica sono interconnesse in modo end-to-end, spesso con
gradienti che fluiscono attraverso entrambe: è il caso di DeepProbLog, dove
le probabilità di regole logiche sono parametri appresi dalla rete
neurale~\cite{deraedt2019deepproblog}, oppure di sistemi dove i vincoli
simbolici sono incorporati come strati differenziabili.

Nell'ambito clinico, Yu et al. argomentano che l'accoppiamento debole è
spesso preferibile perché preserva la verificabilità indipendente dei due
moduli~\cite{yu2023neuro}: un auditor può ispezionare le regole simboliche
senza dover comprendere i pesi della rete neurale, e viceversa. Questa
proprietà è particolarmente preziosa in contesti regolamentati come la
rimborsabilità farmaceutica, dove il medico prescrittore deve poter
giustificare la decisione di fronte a un controllo normativo senza fare
riferimento a meccanismi opachi. Il sistema descritto in questa tesi adotta
un accoppiamento debole con separazione architettuale esplicita: la componente
simbolica, il rule engine che decide la rimborsabilità, non riceve
mai output dalla componente neurale né viene influenzata da essa; la
componente neurale, la pipeline RAG che produce la spiegazione in
linguaggio naturale, riceve invece i risultati della componente simbolica
come contesto strutturato da citare e amplificare.

La separazione non è soltanto una scelta di ingegneria del software: è una
garanzia di correttezza formale. La decisione di rimborsabilità è una
funzione deterministica dei dati clinici del paziente e del testo della Nota
AIFA; non dipende dal modello linguistico, dalla sua versione, dal suo
sampling o dal suo stato interno. Questo è esattamente ciò che differenzia
il sistema sia dai CDSS rule-based puri, che non producono spiegazione in
linguaggio naturale, sia dai sistemi RAG clinici, che non garantiscono
la correttezza normativa della decisione. I modelli basati su grandi modelli
linguistici come Med-PaLM~\cite{singhal2023medpalm} o i sistemi LLM per la
medicina in generale~\cite{thirunavukarasu2023llm} ottengono risultati
impressionanti su benchmark clinici, ma la loro correttezza rimane
probabilistica e non auditabile riga per riga; non sono adatti a sostituire
un sistema normativo dove ogni errore ha conseguenze legali e sanitarie.

La separazione fra decisione deterministica ed esplicazione probabilistica
è pertanto non una limitazione del sistema ma il suo principale vantaggio
architetturale: essa rende la componente di decisione formalmente verificabile
a prescindere dalla qualità della generazione testuale, e rende la componente
di generazione migliorabile o sostituibile senza toccare la logica normativa.


% ------------------------------------------------------------------
\section{Valutazione di sistemi RAG e CDSS}
\label{sec:metriche-rag}
% ------------------------------------------------------------------

Valutare un sistema RAG in un contesto clinico non è riducibile alla misura
di un'unica cifra. Il sistema deve essere corretto nella decisione, preciso
nel retrieval, fedele alle fonti nella generazione e utile all'utente nel
complesso. Questi quattro aspetti sono parzialmente indipendenti: un sistema
può avere un retrieval perfetto e un LLM allucinante, oppure una generazione
fluente che cita fonti inesistenti, oppure, come nel sistema di questa tesi,
una decisione formalmente corretta e una spiegazione di qualità variabile.
Le metriche che seguono, e che sono utilizzate nella valutazione quantitativa
del \Cref{ch:valutazione}, devono essere intese non come giudici assoluti ma
come sonde su dimensioni distinte del comportamento del sistema.

% ──────────────────────────────────────────────────────
\subsection{Metriche di retrieval}
\label{subsec:metr-retrieval}
% ──────────────────────────────────────────────────────

La \emph{Recall@k} misura la frazione di documenti rilevanti effettivamente
restituiti fra i $k$ risultati:
\[
  \text{Recall@}k = \frac{|D_k \cap \text{rel}(q)|}{|\text{rel}(q)|}
\]
dove $D_k$ è l'insieme dei chunk recuperati e $\text{rel}(q)$ l'insieme dei
chunk di riferimento per la query $q$. Con $k$ grande la recall tende a crescere ma la
precisione decade, perché il retriever include sempre più chunk non pertinenti
che occupano la finestra di contesto del modello. La \emph{Precision@k}
complementa la recall misurando la frazione di chunk recuperati che sono
effettivamente rilevanti:
\[
  \text{Precision@}k = \frac{|D_k \cap \text{rel}(q)|}{k}
\]
Nel sistema di questa tesi, la Precision@3 e la Precision@5 sono
particolarmente alte ($0{,}989$ e $0{,}970$ rispettivamente), il che indica
che quasi tutti i chunk recuperati sono pertinenti. Questo risultato è
spiegabile con la struttura delle Note AIFA, dove ogni chunk corrisponde a
un criterio autonomo e i criteri non pertinenti sono raramente recuperati
perché condividono poco vocabolario con la query del caso in esame.

Il \emph{Mean Reciprocal Rank} (MRR) valuta la posizione del primo documento
rilevante nella lista ordinata, mediata su tutte le query:
\[
  \text{MRR} = \frac{1}{|Q|} \sum_{q \in Q} \frac{1}{\text{rank}_1(q)}
\]
dove $\text{rank}_1(q)$ è la posizione del primo chunk rilevante per $q$.
Un MRR pari a $1{,}0$ indica che il chunk rilevante è sempre in prima
posizione: è il risultato osservato nel sistema di questa tesi, coerente
con la struttura di retrieval two-stage dove il cross-encoder affina
l'ordinamento del bi-encoder. Il \emph{Normalized Discounted Cumulative Gain}
(NDCG@k) generalizza MRR pesando la posizione di tutti i documenti rilevanti
con una funzione logaritmica decrescente, ma risulta meno informativo in
questo contesto perché il corpus ha una struttura binaria (un chunk è
rilevante o non lo è) e la Recall@k ne cattura già il comportamento saliente.

% ──────────────────────────────────────────────────────
\subsection{Metriche di fedeltà}
\label{subsec:metr-faithfulness}
% ──────────────────────────────────────────────────────

La \emph{fedeltà} (\emph{faithfulness}) di una risposta misura in quale misura
le affermazioni contenute nella risposta sono supportate dai chunk recuperati,
indipendentemente dalla loro correttezza assoluta rispetto alla realtà. Una
risposta può essere fattualmente corretta ma non fedele ai chunk (se il
modello attinge alla propria memoria parametrica) o fedele ai chunk ma
fattualmente errata (se i chunk stessi contengono un errore). In un contesto
normativo, la fedeltà è la proprietà più critica, perché garantisce che la
spiegazione prodotta dal sistema sia ancorata ai documenti ufficiali e non
a inferenze autonome del modello.

Il framework RAGAS~\cite{ragas2024} implementa la fedeltà come stima della
frazione di affermazioni nella risposta entailed dai chunk di contesto: il
modello LLM-giudice decompone prima la risposta in proposizioni atomiche,
poi verifica per ciascuna se è logicamente implicata dal contesto recuperato,
e infine calcola la proporzione delle proposizioni supportate. Questo approccio
è semantico ma dipende dalla qualità dell'LLM-giudice e risente dell'ambiguità
nelle definizioni di entailment, motivo per cui la versione \emph{strict},
che richiede entailment letterale anziché semantico, è computata
separatamente come limite superiore di rigore.

Un approccio complementare si basa sull'inferenza del linguaggio naturale
(\emph{Natural Language Inference}, NLI): un modello NLI multi-label (nel sistema di questa tesi, a seconda
della metrica, il checkpoint \texttt{mDeBERTa-v3-base-mnli-xnli} o la sua
evoluzione \nlimodel{}) classifica ogni frase della risposta come entailed, neutral
o contraddittoria rispetto alla concatenazione dei chunk di contesto, e la
fedeltà NLI è la proporzione di frasi entailed. L'NLI ha il vantaggio di
essere locale e non dipendente dall'LLM-giudice, ma è più rigida della
fedeltà semantica e tende a sottostimare la fedeltà quando la risposta
parafrasa il testo normativo invece di citarlo letteralmente.

% ──────────────────────────────────────────────────────
\subsection{Qualità della risposta, compositi e statistica}
\label{subsec:metr-quality}
% ──────────────────────────────────────────────────────

La fedeltà valuta la relazione fra risposta e contesto, ma non misura se
la risposta sia utile per chi ha posto la domanda. La \emph{Answer Relevancy}
di RAGAS~\cite{ragas2024} misura invece la coerenza fra la risposta e la
domanda originale: il modello genera $n$ domande sintetiche a partire dalla
risposta prodotta, le codifica come vettori di embedding, e calcola la
similarità coseno media fra i vettori delle domande sintetiche e il vettore
della domanda originale. Un punteggio alto indica che la risposta contiene
l'informazione necessaria a rispondere alla domanda; uno basso segnala
che la risposta è fuori tema o generica.

QAEval~\cite{deutsch2021qaeval} offre un approccio alternativo alla valutazione
della qualità basato su micro-domande: a partire dalla risposta di
riferimento (\emph{gold answer}), viene generato un insieme di domande
di verifica elementari, ciascuna delle quali sonda un fatto specifico
della risposta attesa. La risposta candidata è poi valutata controllando
quante di queste micro-domande ricevono la risposta corretta, producendo
una misura di copertura fattuale fine-grained. Il QAEval F1 combina
precision e recall sulle micro-domande analogamente a una misura di
information extraction. Nel sistema di questa tesi il QAEval F1 medio
è $0{,}172$, un valore basso che riflette la difficoltà di coprire
tutti i dettagli normativi della risposta di riferimento con un modello locale di
8 miliardi di parametri quantizzato a 4 bit (\llmmodel).

Le metriche descritte finora sono eterogenee per scala, natura e peso
diagnostico. Per produrre una valutazione globale del sistema occorre
aggregarle in modo esplicito e giustificato. Il composito \emph{ALES}
(\emph{Aggregated LLM Evaluation Score}) adottato in questa tesi è una
media pesata di quattro componenti:
\[
  \text{ALES} = \alesalpha \cdot D + \alesbeta \cdot E + \alesgamma \cdot U + \alesdelta \cdot Q
\]
dove $D$ è il \emph{DecisionScore} (correttezza della decisione normativa),
$E$ è l'\emph{EvidenceSupport} (fedeltà alle fonti, combinazione pesata con pesi
di base $0{,}50$ per la copertura \emph{verbatim} a 3-grammi (il componente
dominante), $0{,}20$ per l'entailment NLI e $0{,}30$ per la RAGAS
faithfulness; pesi che già sommano a 1 e sono quindi quelli realmente usati
nel calcolo, poiché la quarta componente prevista, la variante \emph{strict}
della RAGAS faithfulness (che ridistribuirebbe i pesi a $0{,}435$/$0{,}174$/
$0{,}261$/$0{,}130$), non è prodotta dal backend RAGAS impiegato e non
contribuisce mai al punteggio), $U$ è la \emph{ContextualUtility}
(media pesata $0{,}75$ della \emph{context precision} e $0{,}25$ della
\emph{context recall}, poiché la \emph{context relevancy} non è restituita dal
backend RAGAS adottato) e $Q$ è la \emph{AnswerQuality} (combinazione $0{,}70$ di
QAEval F1 e $0{,}30$ di Answer Relevancy). Il peso dominante assegnato a $D$ ($\alesalpha$) riflette
la gerarchia di priorità del sistema: la correttezza della decisione normativa
è il requisito non negoziabile, e gli altri componenti ne misurano la
qualità della spiegazione. I valori dei pesi sono impostati come costanti nel
modulo di valutazione e non vengono ottimizzati sui dati di test, per evitare
di adattare la misura ai risultati.

Poiché molte metriche producono proporzioni o medie su campioni di dimensione
finita, la loro incertezza statistica deve essere quantificata per non
confondere la variabilità campionaria con differenze sistemiche. Il metodo
del \emph{bootstrap percentile}~\cite{efron1979bootstrap} ricampiona con
rimpiazzo l'insieme dei $n$ casi per $B = 5000$ iterazioni e calcola
l'intervallo di confidenza al 95\% come il 2,5° e il 97,5° percentile della
distribuzione dei campioni bootstrap. Per le proporzioni binarie (come il
\emph{pass rate} del rule engine) è disponibile anche l'intervallo esatto
di Wilson~\cite{wilson1927ci}, che per proporzioni vicine a 0 o 1 è
superiore al bootstrap percentile perché non richiede approssimazione normale.
L'intervallo di Wilson per una proporzione osservata $\hat{p}$ su $n$
osservazioni con livello di confidenza corrispondente a $z$ è:
\[
  \frac{\hat{p} + \frac{z^2}{2n} \pm z\sqrt{\frac{\hat{p}(1-\hat{p})}{n} + \frac{z^2}{4n^2}}}{1 + \frac{z^2}{n}}
\]
Con $\hat{p} = 1{,}000$ e $n = \goldn$ casi, l'intervallo di Wilson al 95\%
è $[0{,}970;\,1{,}000]$: il rule engine ottiene una classificazione senza errori
sul corpus di riferimento, con un limite inferiore conservativo del 97\%.


I tre pilastri concettuali (reasoning simbolico con logica trivalente,
pipeline RAG con retrieval two-stage, integrazione neuro-simbolica a basso
accoppiamento) non sono stati scelti per completezza enciclopedica ma
perché ciascuno risponde a un requisito preciso del sistema descritto nel
\Cref{ch:progettazione}, dove i loro contorni concettuali diventano contorni
architetturali e i loro vincoli formali diventano vincoli di sicurezza
clinica.

È anche il vocabolario di queste pagine a rendere formulabili le domande di
ricerca. La logica trivalente fissa il significato delle tre decisioni che RQ1
mette alla prova; le metriche di retrieval (Recall@$k$, MRR) sono le sonde di
RQ2; le metriche di fedeltà e il composito ALES traducono in numeri le domande
sulla spiegazione (RQ3 e RQ6) e sul giudizio complessivo (RQ7); gli intervalli
di confidenza bootstrap e di Wilson danno alle risposte la cautela statistica
che ciascuna richiede. Gli strumenti, insomma, non valgono per sé, ma per le
domande che permettono di porre in modo misurabile.
